% BidPilot (GigBid) - R&D report. Compile: latexmk -pdf report.tex  (or upload the whole report/ folder to Overleaf).
% STANDALONE copy, built by scripts/build_report_tables.py. Tables are generated from evals/results/*.json by `python -m scripts.build_report_tables` into report/generated/.
\documentclass[10pt,a4paper]{article}
\usepackage[margin=1.9cm]{geometry}
\usepackage[T1]{fontenc}
\usepackage{lmodern,microtype,booktabs,xcolor,enumitem,graphicx,caption,float}
\usepackage[hidelinks]{hyperref}
\setlist{nosep,leftmargin=*}
\captionsetup{font=small,labelfont=bf}
\setlength{\textfloatsep}{8pt plus 2pt minus 2pt}
\newcommand{\pending}{\textcolor{orange!80!black}{\textit{not measured}}}
\newcommand{\file}[1]{\texttt{\detokenize{#1}}}
\newcommand{\pendingtable}[2]{\begin{center}\fbox{\parbox{0.9\linewidth}{\small\textbf{#1:} \pending{} at this commit; source file \file{#2}.}}\end{center}}
\emergencystretch=3em
\newcommand{\ragvzeroRfive}{0.544}
\newcommand{\ragvzeroNdcg}{0.584}
\newcommand{\ragvoneRfive}{0.557}
\newcommand{\ragvoneNdcg}{0.587}
\newcommand{\ragvtwoRfive}{0.539}
\newcommand{\ragvtwoNdcg}{0.581}
\newcommand{\decF}{0.645}
\newcommand{\decScamP}{0.667}
\newcommand{\decScamR}{1.000}
\newcommand{\decInj}{1.000}
\newcommand{\decPrice}{1.000}
\newcommand{\decPninetyfive}{44.42}

\providecommand{\decF}{\pending}\providecommand{\decScamP}{\pending}\providecommand{\decScamR}{\pending}
\providecommand{\decInj}{\pending}\providecommand{\decPrice}{\pending}\providecommand{\decPninetyfive}{\pending}

\title{\textbf{BidPilot: a multi-agent freelancer copilot}\\[2pt]\large R\&D report (AI-4016, GigBid) -- engineering design review}
\author{Qasim Hafeez (22F-3812)\\\small Code: \url{https://github.com/qasim-hafeez00/bidpilot} \quad Deployment: \url{https://bidpilot-e409.onrender.com}}
\date{November 2026}

\begin{document}
\maketitle
\begin{abstract}\noindent
BidPilot implements GigBid: LangGraph agents find freelance jobs, decide whether a job is worth a bid, write proposals
grounded in portfolio evidence and handle client conversations, while a human approves every proposal, message and
Connects spend. All marketplace access goes through one port with a GigSim adapter and a live Upwork MCP adapter, and every
agent reaches its tools only through MCP. The system runs on a CPU laptop with open-weight models (Ollama, Qwen2.5 3B/7B)
and on hosted open-weight models (gpt-oss-20b/120b via Groq). The main findings: (1) the pre-registered
retrieval hypothesis was \emph{refuted}: requirement decomposition and hybrid BM25 did not beat naive dense retrieval on
multi-requirement jobs, and only multi-vector project aggregation helped; (2) claim verification plus an independent fact checker
keeps unsupported claims to a fraction of one per proposal, but the checker found cases the in-loop verifier missed;
(3) an LLM judge agreed moderately with reference labels (quadratic $\kappa = 0.53$, $n=28$); (4) the hand-built graph used about half
the tokens of a ReAct agent or a supervisor on the same jobs; (5) the X2 A/B test favoured problem-first openings (25\% vs.\ 6.7\% replies),
but the difference was not significant ($p=0.33$). All numbers come from files in
\file{evals/results/}. Sample sizes are small and are stated next to every result.
\end{abstract}

\section{Problem, constraints and budgets}
\textbf{Workflow.} A freelancer scans feeds, spends paid Connects, writes proposals and handles replies. BidPilot automates
discovery (Scout), assessment (Analyst, Risk, Fit, Pricing) and drafting (Writer, Critic, Compliance, Liaison). The
freelancer keeps every decision that leaves the system.

\textbf{Never automated (enforced in code, not in prompts).} Submitting a proposal, sending a message and spending Connects
all pass through a LangGraph \texttt{interrupt()} (\file{agent/nodes/approval.py}, \file{agent/nodes/client_liaison.py}).
Per-agent MCP allowlists mean that only the post-approval \texttt{submitter} can see \texttt{submit\_proposal}
(\file{agent/mcp_client.py}). The optional autonomy setting can only auto-\emph{skip} low-fit jobs. Live mode only creates
Upwork drafts, which the user confirms on Upwork.

\textbf{Declared budgets} (\file{agent/config.yaml}): p95 latency $\le$ 20\,s per job analysed and $\le$ 300\,s per proposal;
cost $\le$ \$0.002 per job analysed and $\le$ \$0.02 per proposal on the hosted profile (\$0 on the local profile); graph budgets of 30
steps, 16k tokens and at most 2 Critic revisions per thread.

\textbf{Compute reality.} Development ran on an Intel i7-8550U (4 cores, 16\,GB, no GPU). Measured Ollama throughput: Qwen2.5-3B
6.1 tok/s and Qwen2.5-7B 2.9 tok/s. This shaped the design: rules-first nodes, LLM calls only where they add value, an LLM
cache, and hosted open-weight models (\texttt{open\_hosted} profile) for the LLM-heavy experiments, which would otherwise have taken days on this CPU.

\section{Data}
The instructor's GigSim, persona corpora and dev set were not in the starter kit I received, so I rebuilt them
(\file{scripts/build_data/}), keeping real and synthetic data distinct. Every record carries a \texttt{provenance} field; full
details, licences and limitations are in \file{docs/DATA_CARD.md}.
\begin{itemize}
\item \textbf{Real:} 27{,}713 usable Upwork postings with skill tags, client statistics and Connects cost (Kaggle
  \emph{upwork-freelance-jobs-60k}); pricing percentiles per domain and rate type from the 53k-posting
  \emph{upwork-job-postings-dataset-2024}; scam phrasing from the labelled fraudulent postings in \emph{trial-1}.
\item \textbf{Synthetic:} three personas (AI/ML, full-stack, data analyst), each with a PDF CV, 17--18 case studies, repositories
  (README plus code), reviews and 150 past proposals labelled won/replied/ignored. The corpus is built to defeat naive RAG:
  \emph{implied} skills appear only in code imports (\texttt{from celery import shared\_task}) or paraphrased prose (``a background
  worker queue''), and each project is described three ways. Client hidden preferences, scam payloads, injected instructions and
  screening questions are also synthetic.
\item \textbf{Rejected:} \emph{freelancer-success} (hourly rates capped; median \$280 for every skill) and the 1.3M-row contracts
  dataset (redundant with the postings for pricing).
\item \textbf{Splits:} a dev set of 200 jobs per persona with gold bid/skip, scam and injection labels, graded gold evidence and a fair
  price band; a git-ignored held-out split of 80 jobs per persona with unseen clients stands in for the hidden test set.
\end{itemize}
\textbf{Threat to validity.} I wrote both the persona generator and the system, so they share vocabulary. Treat the results as
an upper bound on what a stranger's portfolio would give.

\section{Architecture}
\begin{figure}[H]\centering\scriptsize
\begin{verbatim}
 Scout (worker, cron) --job_id--> intake --Send--> Job Analyst ----\
                                         --Send--> Risk Assessor ---> decide --skip/watch--> END
                                         --Send--> Fit Matcher -----/     \--bid--> Pricing -> Writer -> verify claims
   Critic --fail (<=2x)--> Writer;  Critic --pass--> Compliance --block--> END (reason shown in UI)
   Compliance --> dry-run preview --> [interrupt: approve | edit | reject] --> submit (idempotent) --> END
 Conversation graph: read messages -> classify intent -> (find slots) -> draft reply -> [interrupt] -> send (+ tentative hold)
\end{verbatim}
\caption{Both graphs are hand-written \texttt{StateGraph}s (\file{agent/graphs/}); Mermaid renderings are in \file{docs/images/}.}
\end{figure}
\textbf{Why each agent exists.} The Analyst produces a typed \texttt{JobSpec} (rules first, small LLM for hidden needs). The Risk
Assessor needs the client-history tool and a lexicon calibrated on scam text. The Fit Matcher is the only reader of the
portfolio and maps each decomposed requirement to evidence. Pricing combines market percentiles, the client's range and the
user's floor with a Connects expected-value check. The Writer is the only node that needs the strongest model. The Critic and
the Compliance Guard are separate because they fail for different reasons (quality vs.\ policy), and only Compliance may block.
The Liaison owns conversations.

\textbf{State, trace and memory are kept apart} (\file{agent/state.py}). \emph{State} holds typed fields with reducers wherever
parallel branches write, including a severity-ordered \texttt{status} reducer. The \emph{trace} is append-only records plus
LangSmith/Phoenix spans plus the MCP audit log, linked by one trace id from the UI request to the audit row. \emph{Long-term memory}
is a LangGraph store with two namespaces, written only through \file{agent/memory/promotion_rules.py}: a style rule is promoted
when an edit pattern appears in $\ge$3 of the user's last 5 edits (demoted when it appears in none), and an outcome insight
only when both arms have $n\ge5$ and their reply rates differ by $\ge$15 points. Hand-offs are Pydantic models
(\file{agent/handoffs.py}).

\textbf{MCP surface} (\file{docs/TOOLS.md}). Three servers built on the official Python SDK, each with Pydantic input schemas,
declared output schemas and \texttt{readOnly}/\texttt{destructive}/\texttt{idempotent} annotations: \emph{Portfolio} (read-only
search, case studies, skills; resources and prompts), \emph{Marketplace} (8 read and 3 write tools, required idempotency keys,
\texttt{dry\_run}, per-user JWT bearer on Streamable HTTP, a per-user rate limiter and an append-only audit log) and \emph{Workspace}
(free slots, tentative holds, email drafts on a documented SQLite mock). The third-party servers are the official Upwork MCP
server (OAuth~2.1 + PKCE, tokens Fernet-encrypted per user) and the GitHub MCP server for repository ingestion.

\textbf{Marketplace port (R1).} \file{agent/marketplace_port.py} defines the interface; \texttt{GigSimAdapter} (REST, bounded
retries honouring \texttt{Retry-After}) and \texttt{UpworkLiveAdapter} (\texttt{MultiServerMCPClient}, with its tool mapping in config
and validated against \texttt{list\_tools}) implement it. Switching \texttt{marketplace.mode} changes only which adapter the
Marketplace server instantiates.

\textbf{Durability and failure handling.} Postgres holds the checkpointer and the store. Approvals resume the interrupted thread
from the UI, and the pending payload is read from the checkpoint, so a refresh never loses it. Submit idempotency keys are a pure
function of (user, job, draft version). Every node runs under a guard (\file{agent/graphs/_guard.py}) that enforces step and token
budgets, bounded retries with timeouts, a \emph{degraded mode} (\texttt{MarketplaceUnavailable} sets \texttt{queued\_degraded},
the UI shows a banner and the worker retries) and a dead-letter status. Background work runs on Arq/Redis with deterministic
job ids, and a reconciler re-enqueues runs that were in flight when a worker died (an integration test kills a run mid-graph and
checks exactly-once submission on resume).

\section{R\&D question B -- retrieval, claim verification and what ships}\label{sec:rag}
\textbf{Hypothesis (stated before running).} Requirement decomposition plus hybrid retrieval raises evidence Recall@5 by
$\ge$10 points over V0 on jobs with four or more required skills, because implied skills live in code chunks that whole-post
dense retrieval misses.

\textbf{Setup.} Every variant is a preset of one composable pipeline (\file{rag/pipeline.py}), so the ablation removes one flag at
a time. Embeddings: bge-small-en-v1.5; cross-encoder: ms-marco-MiniLM-L6 (both ONNX on CPU). Gold relevance of a project is its
overlap with the job's skills (explicit and implied), with ubiquitous skills (Python, SQL, \ldots) weighted 0.25. Two runs exist: all
seven variants on a 128-job sample (first 60 dev jobs per persona, minus scams; Table~\ref{tab:rag60}, from the run log
\file{evals/results/rag_limit60_run.log}), and V0--V2 on the full dev set (396 jobs, Table~\ref{tab:rag}). The composite and V8 runs did not
finish before submission (a shared SQLite cache lock crashed the run), so they are not reported.

\begin{table}[t]\centering\small
\caption{All variants on the 128-job sample (62 with $\ge$4 skills). Source: \file{evals/results/rag_limit60_run.log} of the \texttt{--limit 60} run; its JSON was later
overwritten by the full-dev run, so this table is transcribed from the log.}\label{tab:rag60}
\begin{tabular}{lrrrrr}\toprule Variant & R@5 & MRR & nDCG@10 & R@5 ($\ge$4 skills) & p95 s\\\midrule
V0 naive dense & 0.528 & 0.704 & 0.595 & 0.554 & 0.00\
V1 decomposition & 0.547 & 0.670 & 0.576 & 0.537 & 0.01\
V2 hybrid BM25+dense & 0.523 & 0.669 & 0.571 & 0.515 & 0.00\
V3 + cross-encoder & 0.533 & 0.642 & 0.552 & 0.530 & 16.70\
V4 outcome-aware & 0.521 & 0.626 & 0.554 & 0.497 & 0.01\
V5 multi-vector & 0.559 & 0.661 & 0.580 & 0.564 & 0.00\
V6 contextual & 0.539 & 0.653 & 0.564 & 0.533 & 0.00\
\bottomrule\end{tabular}\end{table}

\begin{table}[t]\centering\small
\caption{Retrieval quality on the full dev set (396 non-scam jobs with gold evidence, 3 personas). Source: \texttt{evals/results/rag\_variants.json}.}\label{tab:rag}
\begin{tabular}{lrrrrr}\toprule Variant & R@5 & MRR & nDCG@10 & $\Delta$R@5 (pts) & p95 s\\\midrule
V0 naive dense & 0.544 & 0.661 & 0.584 & 0.0 & 0.03 \\
V1 decomposition & 0.557 & 0.653 & 0.587 & +1.3 & 0.09 \\
V2 hybrid BM25+dense & 0.539 & 0.662 & 0.581 & -0.5 & 0.07 \\
\bottomrule\end{tabular}\end{table}

\textbf{Result: the hypothesis is refuted.} On the 62 multi-requirement jobs, V0 reached Recall@5 0.554. Decomposition (V1) reached
0.537 ($-1.7$ points) and hybrid BM25+dense (V2) 0.515 ($-4.0$); the cross-encoder (V3) did not recover the gap (0.530). The
only consistent gain came from \emph{multi-vector project aggregation} (V5: 0.564 on multi-requirement jobs, 0.559 overall,
$+3.1$ points), which scores a project by its best chunks across case study, README, code and proposal. The 95\% intervals
at $n=128$ are roughly $\pm$5 points, so none of these differences is large. The mechanism matches the earlier smoke run: BM25
over code rewards files whose identifiers are common words, and decomposed single-skill queries match generic chunks
(``Python'') as strongly as the specific ones. MRR is highest for V0 (0.704), so the naive retriever usually gets the top hit
right and loses on coverage. The full dev set (396 jobs, 163 with $\ge$4 skills) confirms the direction: on multi-requirement jobs
V0 0.560, V1 0.550, V2 0.532.

\begin{table}[t]\centering\small\caption{Recall@5 by persona and job type. Source: \texttt{rag\_variants.json}.}\label{tab:ragbreak}
\begin{tabular}{lrr}\toprule Slice & V0 naive dense & V2 hybrid BM25+dense\\\midrule
ai\_ml & 0.574 & 0.577 \\
fullstack & 0.556 & 0.549 \\
data\_analyst & 0.504 & 0.493 \\
\midrule
adjacent & 0.418 & 0.359 \\
few\_req (<4 skills) & 0.533 & 0.545 \\
in\_domain & 0.557 & 0.561 \\
injection & 0.591 & 0.577 \\
multi\_req (>=4 skills) & 0.560 & 0.532 \\
off\_domain & 0.317 & 0.237 \\
\bottomrule\end{tabular}\end{table}
\pendingtable{Ablation of the composite pipeline}{evals/results/ablation.json}

\textbf{Claim verification (B3)} (\file{rag/verify.py}). Each first-person experience or result sentence in a draft is a claim,
labelled supported, partly supported or unsupported by the cross-encoder score against the cited (else retrieved) evidence, plus
entity checks: every skill and number in the claim must appear in the evidence. Unsupported claims are removed before the Critic
sees the draft. The unsupported-claim \emph{rate} is measured separately, by a stricter checker built on the persona's
ground-truth facts (\file{evals/checkers/fact_check.py}), so the system is not graded by its own verifier.

\begin{table}[t]\centering\small\caption{Proposal-level results (open\_hosted profile; 5 gold-bid jobs per persona). Unsupported claims are counted by the independent fact checker; reply rate is from GigSim's rule-based clients. Source: \texttt{proposals.json}, \texttt{judge.json}.}\label{tab:prop}
\begin{tabular}{lrrrrrrrrr}\toprule Retrieval & n & Unsup.\ before & Unsup.\ after & Critic loops & Judge (1--5) & Reply rate & p95 s & Tokens & USD\\\midrule
Composite (shipped) & 15 & 0.20 & 0.40 & 2.80 & 3.21 & 0.14 & 286 & 4888 & 0.0012 \\
V0 naive dense & 15 & 0.27 & 0.20 & 2.80 & 3.14 & 0.14 & 168 & 4783 & 0.0010 \\
\bottomrule\end{tabular}\end{table}

\textbf{Proposal-level results} (Table~\ref{tab:prop}; 15 gold-bid jobs per configuration, Writer gpt-oss-120b, small nodes
gpt-oss-20b). The verifier removed 1.0--1.3 claims per proposal. Even so, the independent checker still found 0.2--0.4
unsupported claims per final proposal, and for the shipped retriever the final rate (0.40) was \emph{higher} than the raw draft's
(0.20). The cause was the Critic loop (below): each rejected draft was rewritten, and a rewrite can add claims that the
cross-encoder rates as ``partly supported'' while the fact checker does not. Judge scores (3.21 vs.\ 3.14) and GigSim reply
rates (0.14 vs.\ 0.14, $n=14$ submitted each) do not separate the two retrievers; at this sample size a difference below about
25 points in reply rate cannot be detected. Cost was \$0.0010--0.0012 per proposal (budget \$0.02). p95 latency was 168\,s for V0
and 286\,s for the shipped pipeline (budget 300\,s), measured while other evaluations shared the CPU.

\textbf{Judge calibration.} The judge (Qwen3, a different model family from the writer, to limit self-preference) scored 28 proposals.
These were compared with reference labels on the same 1--5 rubric, given blind to configuration. \noindent\textbf{Judge calibration.} calibrated: quadratic-weighted $\kappa$ = 0.53 on $n$ = 28 labelled proposals (1 disagree by $\ge$2 points). Source: \texttt{judge.json}.
\emph{Disclosure:} the reference labels in \file{evals/datasets/judge_labels.csv} were written by an AI assistant (Claude) at my
request, not by me. $\kappa$ therefore measures agreement between two models, not between the judge and a human. The judge
gave 4 to half the proposals and rarely used 1 or 5, so its scores are compressed. The blind classmate panel was not run.

\textbf{Shipped choice.} \texttt{rag.variant} is the composite pipeline (multi-vector aggregation over hybrid candidates with
reranking), chosen because V5's aggregation is the only component with a measured gain. The composite itself was not measured
end to end on the retrieval set, so the honest recommendation from this data is V5's preset without hybrid or decomposition. The \texttt{search\_evidence} tool disables query
decomposition, since the Fit Matcher already decomposes at the agent level and doing it twice cost 16--25\,s per call. On the
512\,MB free hosting tier the cross-encoder is replaced by a bi-encoder score (\texttt{RERANK\_DISABLED=1}), which keeps the app
within memory (measured 276\,MB vs.\ 373\,MB resident).

\section{R\&D question A -- tool granularity}
\textbf{Hypothesis.} Fewer, coarser tools reduce tool-selection errors for a small model, while rewritten descriptions recover
argument accuracy. \textbf{Setup} (\file{evals/runners/tool_granularity.py}): 40 tasks with a known correct call, one turn, three
tool sets bound to the same model: the 14 real tools with their terse descriptions, 3 coarse tools with an \texttt{action} enum,
and the 14 tools with rewritten ``use when\ldots'' descriptions.
\begin{table}[t]\centering\small\caption{R\&D A: tool granularity with \texttt{openai/gpt-oss-20b}. Source: \texttt{tool\_granularity.json}.}\label{tab:tools}
\begin{tabular}{lrrrrr}\toprule Tool set & n & Tool sel.\ acc. & Arg.\ acc. & Tokens/task & Latency s\\\midrule
Fine-grained (14 tools) & 40 & 1.000 & 1.000 & 508 & 2.11 \\
Coarse (3 tools + action) & 40 & 0.850 & 0.300 & 296 & 2.24 \\
Fine + rewritten descriptions & 40 & 0.975 & 0.975 & 759 & 5.02 \\
\bottomrule\end{tabular}\end{table}
\textbf{Result: the hypothesis is refuted for this model.} With gpt-oss-20b, the fine-grained set was already at ceiling (40/40 correct
tool and arguments). Collapsing to three coarse tools cut tokens per task by 42\% (508 to 296) but lowered tool selection to 0.85
and argument accuracy to 0.30. With a free-form \texttt{args} object, the schema no longer tells the model which fields an action
needs, so the expected values often went missing (per-call logs were not kept, so the failure modes are not broken down further). Rewritten descriptions added 49\% tokens (759) and gave 0.975 on both, one error more than the terse set. \textbf{Decision:}
keep fine-grained tools with short descriptions; coarse tools trade a small token saving for a large loss in argument
accuracy. \emph{Caveats:} the 40 tasks are single-turn and partly paraphrased duplicates, so a 20B model saturates them; the original
target, Qwen2.5-3B, was not run on this set because of the CPU time it would take, and a small model may well behave differently.

\section{R\&D question C -- is the multi-agent design worth it?}
\textbf{Setup} (\file{evals/runners/arch_compare.py}): the same jobs, model and in-memory MCP tools run through (a) one ReAct agent
holding all tools, (b) a \texttt{langgraph-supervisor} with research, portfolio and writer workers, and (c) the hand-built graph.
\begin{table}[t]\centering\small\caption{R\&D C: the same jobs (3 per persona) through three architectures. Source: \texttt{arch\_compare.json}.}\label{tab:arch}
\begin{tabular}{lrrrrrrrr}\toprule Architecture & n & Bid/skip acc. & Scams bid & Unsup./prop. & Policy viol. & Tokens/job & p95 s & Errors\\\midrule
Single ReAct agent & 9 & 0.556 & 0 & 0.00 & 1 & 3936 & 50 & 0 \\
LangGraph supervisor & 9 & 0.667 & 0 & 1.33 & 0 & 3731 & 30 & 0 \\
Hand-built graph & 9 & 0.667 & 0 & 0.33 & 0 & 1702 & 114 & 0 \\
\bottomrule\end{tabular}\end{table}
\textbf{Result} (9 jobs, 3 per persona, gpt-oss on Groq). The hand-built graph and the supervisor both decided 6 of 9 jobs correctly; the
single ReAct agent decided 5. The graph used 1{,}702 tokens per job against 3{,}731 (supervisor) and 3{,}936 (ReAct), because its rule-based
nodes skip LLM calls on jobs that should be skipped. The supervisor's proposals had 1.33 unsupported claims each against 0.33 for the graph,
because nothing between its writer and the output checks claims. The ReAct agent left one policy violation in its output; the graph and
the supervisor left none. No architecture bid on a scam. The graph had the worst p95 latency (114\,s), because it runs claim
verification and up to two Critic loops. \textbf{Conclusion:} at $n=9$ the accuracy difference is noise; the measured benefits of the hand-built
graph are half the tokens and far fewer unsupported claims, paid for with latency.

\textbf{Decision quality of the hand-built graph} (rules-only mode, \texttt{LLM\_DISABLED=1}, so decisions are deterministic):
\begin{table}[t]\centering\small\caption{Decisions on the dev set (graph in analysis-only mode through the in-memory MCP client). Source: \texttt{decision\_eval.json}.}\label{tab:dec}
\begin{tabular}{lrrrrrrrrr}\toprule Slice & n & Bid F1 & P & R & Scam P & Scam R & Inj.\ det. & Price in band & p95 s\\\midrule
All & 60 & 0.645 & 0.714 & 0.588 & 0.667 & 1.000 & 1.000 & 1.000 & 44.42 \\
ai\_ml & 20 & 0.444 & 0.667 & 0.333 & 0.500 & 1.000 & 1.000 & 1.000 & 35.36 \\
fullstack & 20 & 0.500 & 0.500 & 0.500 & 1.000 & 1.000 & \pending{} & 1.000 & 37.97 \\
data\_analyst & 20 & 0.778 & 0.778 & 0.778 & 0.500 & 1.000 & 1.000 & 1.000 & 54.87 \\
\bottomrule\end{tabular}\end{table}
On 60 dev jobs: bid F1 0.65 (precision 0.71, recall 0.59), accuracy 0.82, scam recall 1.00 at precision 0.67 (4 legitimate jobs flagged),
all injections detected with no false positives, and every quoted price inside the fair band. Recall is the weak point: 7 gold bids were
skipped, mostly for the AI/ML persona, where the fit score sits just below the 0.55 threshold. The p95 analysis latency of 44\,s
breaks the 20\,s budget; it was measured while three other evaluation processes shared the 4-core CPU, and the cross-encoder dominates.

\section{The change request}
GigSim's 2 November policy update was not available to me, so I \emph{simulated} one and label it as such (\file{gigsim/app.py},
policy v2). It has two parts: (1) proposals may not contain external links (submission returns HTTP 422), and (2) premium jobs (Expert
level or budget $\ge$\$1{,}000) cost 4 extra Connects. The agents needed no code change. The Compliance Guard reads
\texttt{policy.v2\_rules} from config, enabled by \texttt{marketplace.policy\_version: 2}, and Pricing reads the per-job Connects cost
that the marketplace returns. The diff is one line of \file{agent/config.yaml} plus the policy-pack text (\file{docs/CHANGE_REQUEST.md}).
This worked because rules and prompts were configuration from the start. A change that needs new \emph{logic} (e.g.\ a weekly proposal quota)
would still need code in \texttt{decide}.

\section{Signature extension -- X2 proposal experiments}
Two opening strategies (problem-first vs.\ question-first) run as an A/B test in GigSim, analysed with a two-proportion $z$-test and
Fisher's exact test on reply rate (\file{evals/runners/ab_test.py}, \file{docs/EXTENSION.md}). Each submitted proposal's opening style is
recorded and feeds outcome memory, so only a winner that clears the promotion rule becomes an insight that the Writer reads.
\begin{table}[t]\centering\small\caption{X2 A/B test on GigSim reply rate. Two-proportion $z$ = 1.39 ($p$ = 0.166), Fisher exact $p$ = 0.333. Source: \texttt{ab\_test.json}.}\label{tab:ab}
\begin{tabular}{lrrr}\toprule Arm & Submitted & Replies & Reply rate\\\midrule
A: problem-first & 16 & 4 & 0.25 \\
B: question-first & 15 & 1 & 0.07 \\
\bottomrule\end{tabular}\end{table}
\textbf{Result} (6 jobs per persona per arm). Problem-first got 4 replies from 16 proposals (25\%) and question-first 1 from 15 (6.7\%).
Two-proportion $z=1.39$, $p=0.17$; Fisher's exact $p=0.33$. This is \emph{not significant}, and the promotion rule correctly did not promote
it into memory. Detecting an 18-point difference at 80\% power needs roughly 70 proposals per arm. Reply rates come from GigSim's rule-based
clients, so they measure fit to that model, not to real clients.

\section{Failure analysis}
Taxonomy: missed evidence, wrong skip, mispriced, unsupported claim, generic opening, compliance miss, injection, tool error.
Failures found, with the fix and the measured effect where one exists:
\begin{enumerate}
\item \textbf{Unsupported claim / Critic loop (found in the proposal run).} gpt-oss writes citations with full-width brackets (U+3010/U+3011 around \texttt{E:id}), but the Critic
  counted only \texttt{[E:id]}. It therefore rejected almost every draft (2.8 Critic passes per proposal; the cap of 2 revisions was
  hit) and the full-width markers leaked into five final drafts that a client would have seen. Fix: normalise both marker forms in
  the Writer and the verifier, with a regression test. \emph{Effect} (rerun on the same 15 jobs, \file{proposals_fixed.json}): no markers
  leaked, and unsupported claims after verification fell from 0.40 to 0.13 per proposal. Critic passes stayed at 2.8, so the Critic
  also rejects for other reasons and its rubric needs inspection. Part of the rerun was served from the LLM cache, so its cost reads \$0
  and its latency (p95 53\,s) is not comparable.
\item \textbf{Instruction vs.\ attention check.} One real posting says ``Include `LFG' in your proposal to show that you read this.''
  The Writer complied, which a human would also do, and the injection detector did not flag it. The boundary between a client's
  screening instruction and a prompt injection is semantic. I kept this behaviour and wrote the rule into the policy pack (rule 5): a
  client's request about the proposal's \emph{content} may be followed, while instructions addressed to an AI, or asking for contact details,
  payment or off-platform moves, never are. A regex detector cannot draw this line reliably; an LLM classifier on flagged sentences could.
\item \textbf{Generic opening.} Two drafts for the chemistry-classification job opened with ``I am a passionate and hardworking AI/ML Engineer'',
  a phrase the style rules forbid. Both received the lowest reference label (2). The opening-style check belongs in the Critic's rubric.
\item \textbf{Missed evidence (off-domain bids).} With \texttt{force\_bid}, the data-analyst persona drafted for an AutoCAD/Revit job using
  VBA evidence. The decision node had correctly not bid (fit below threshold); the failure is that the harness forced it, and the draft
  shows what the Writer does with weak evidence.
\item \textbf{Wrong skip (risk false positive).} One gold-bid AI job was skipped in both configurations despite \texttt{force\_bid}, because the
  risk lexicon scored it as a likely scam ($\ge$0.5), and force-bid never overrides a scam flag. This is safe by design, but it is a false positive.
\item \textbf{Earlier fixes during development:} sentence-only requirements now weigh 0.5 in the fit score (a proof-of-concept job had
  scored 0.40); the gold labeller requires all title skills (a Solidity job was labelled ``bid'' for the full-stack persona); the adapter
  tolerates non-JSON error bodies when another service answers GigSim's port; MCP sessions get one owner task each (anyio crashed on
  cross-task close); and \texttt{search\_evidence} skips re-decomposition (16--25\,s per call before).
\item \textbf{Hosting (tool error).} The first Render deployment restarted in a loop as soon as the worker started an analysis. Loading both
  ONNX models and three MCP servers in one process exceeded the 512\,MB instance. Fix: in-process MCP transport, single-threaded ONNX,
  no cross-encoder on that tier.
\end{enumerate}

\section{Safety and platform respect}
\textbf{Never does:} submit, message or spend Connects without an approval interrupt; scrape, automate a browser or use session cookies;
claim experience the evidence does not support; follow instructions found inside a job post or client message.
\textbf{Enforcement:} untrusted text is PII-masked and delimited at intake, and spoofed delimiters are neutralised. The Risk Assessor detects
injected instructions and strips those sentences before analysis. The Writer prompt forbids following them, and Compliance removes any
echo (e.g.\ ``BANANA'', guarantees, contact details). Writes require approval, allowlists and idempotency keys. PII masking covers logs,
the audit log and trace exports, and OAuth tokens are encrypted and never logged. Every query is scoped by \texttt{user\_id}, and the audit
table is append-only through a trigger. The injection-detection rate on the dev sample is in Table~\ref{tab:dec}.
\textbf{What still gets through:} paraphrased injections that match no pattern reach the Writer as delimited data, where only the prompt
and the compliance echo check defend. The cross-encoder is a relevance model, not an entailment model, so a fluent but subtly wrong
claim near true evidence can pass as ``partly supported''; the proposal run above shows this happening.

\section{Held-out results against the declared budgets}
The instructor's hidden run cannot be reproduced here. The stand-in is the git-ignored held-out split (new clients), run with
\texttt{--split test}. \emph{That run is invalid}: the dev GigSim instance did not load the held-out clients
(\texttt{GIGSIM\_INCLUDE\_TEST=0}), so every client-history call returned 404 and all 45 jobs ended in \texttt{needs\_attention}. The degraded
path behaved as designed (no job was bid on without client data), but no held-out quality number is reported.
Proposal p95 latency (168--286\,s) and cost (\$0.0010--0.0012) are within the 300\,s and \$0.02 budgets on the hosted profile. Analysis latency
on the CPU laptop exceeded the 20\,s budget whenever the cross-encoder ran under load (V3 p95 16.7\,s for retrieval alone), which is a
second reason the hosted tier drops it.
\textbf{Deployment.} \url{https://bidpilot-e409.onrender.com} (Render Blueprint in \file{render.yaml}: app with an inline worker, GigSim,
Postgres and Key Value, all on the free plan). Free instances sleep after 15 minutes, so the first request can take about a minute. A reviewer
account is seeded at start-up (credentials in the README).

\section{Models: open configurations compared}\label{sec:models}
Three open-weight configurations are defined in \file{agent/config.yaml}: \texttt{open} (Qwen2.5-3B for every node, 7B judge),
\texttt{open\_7b\_writer} (7B Writer) and \texttt{open\_hosted} (gpt-oss-20b small, gpt-oss-120b Writer and Critic, Qwen3 judge, on Groq).
Measured local throughput was 6.1 tok/s (3B) and 2.9 tok/s (7B). At those speeds one proposal with up to two Critic loops takes 5--10
minutes, so the proposal-level experiments ran on \texttt{open\_hosted} and the local profiles were exercised only in development and the CI
smoke test. A controlled quality comparison of the local and hosted profiles on the same jobs was not run. No paid model is used.

\appendix
\section{Reproducing every number}
{\scriptsize
\begin{verbatim}
make data                                                    # Kaggle -> personas, GigSim seed, dev/held-out sets (seeded)
python -m evals.runners.rag_eval --limit 60                  # 128-job table; --limit 0 for the full dev set
LLM_DISABLED=1 python -m evals.runners.decision_eval --limit 20   # decision_eval.json
MODEL_PROFILE=open_hosted python -m evals.runners.gigsim_outcomes --configs best,v0 --per-persona 5
python -m evals.runners.judge_calibration    # needs evals/datasets/judge_labels.csv for kappa
python -m evals.runners.tool_granularity; python -m evals.runners.arch_compare --per-persona 3
python -m evals.runners.ab_test --per-persona 6; python -m scripts.build_report_tables
\end{verbatim}}
\section{AI usage}
See \file{docs/AI_USAGE.md}. An AI coding assistant (Claude Code) wrote most of the code under my direction. At my request it also
produced the judge reference labels, as disclosed in \S\ref{sec:rag}. One error it made that was caught: its first data builder applied
about 500 regexes per job across 110k jobs and ran for over 10 minutes; it was replaced by a substring-prefiltered matcher (97\,s end to end).
\end{document}
